% Source: https://artofproblemsolving.com/community/c3h1606878_number_theory
% Stable identity: aops-topic-1606878-q01
\textbf{(a)} As $d$ runs over the positive divisors of $n$, so does $\dfrac {n}{d}$ (the divisors pair up as $d\leftrightarrow\dfrac {n}{d}$). Hence
\[\sum_{d\mid n}\dfrac{1}{d}=\dfrac{1}{n}\sum_{d\mid n}\dfrac {n}{d}=\dfrac{1}{n}\sum_{e\mid n}e=\dfrac{\sigma(n)}{n},\]
where the middle step puts $e=\dfrac {n}{d}$.

\textbf{(b)} $1000=2^3\cdot5^3$, so $\sigma(1000)=(1+2+4+8)(1+5+25+125)=15\cdot156=2340$, and the sum is
\[\dfrac{2340}{1000}=\boxed{\dfrac{117}{50}}.\]

\textbf{(c)} By (a), the sum is $2$ exactly when $\sigma(n)=2n$, the definition of a perfect number. For $n=28$:
\[1+\dfrac{1}{2}+\dfrac{1}{4}+\dfrac{1}{7}+\dfrac{1}{14}+\dfrac{1}{28}=\dfrac{28+14+7+4+2+1}{28}=\dfrac{56}{28}=2 .\]
